\section{Conclusion}

This project presented a data mining analysis of the integrated Heart Disease dataset from the UCI
Machine Learning Repository, with the goal of predicting the presence of heart disease using a
K-Nearest Neighbors (KNN) classifier. The workflow included data preprocessing, an extensive exploratory
data analysis (EDA), and the implementation of a mixed-distance KNN model capable of handling both
numeric and categorical attributes.

The preprocessing stage ensured that the dataset was suitable for distance-based learning. Missing values
were imputed using median and mode strategies for numerical and categorical features, respectively.
Numeric attributes were then normalized using Min--Max scaling to prevent features with large numerical
ranges from dominating the distance computation. An 80/20 train–test split was used to evaluate the model
on unseen data.

The EDA revealed several clinically meaningful patterns. Among all numeric attributes, \texttt{thalach}
(maximum heart rate achieved) and \texttt{oldpeak} (exercise-induced ST depression) showed the strongest
associations with heart disease. Patients with heart disease generally exhibited lower \texttt{thalach}
values and higher \texttt{oldpeak} values, producing the clearest visual separation between classes in the
scatter plots and pair plots. Other variables such as \texttt{chol} and \texttt{trestbps} exhibited weaker
or more diffuse relationships and contained several medically plausible outliers. The target distribution
was nearly balanced, which is favorable for classification without additional rebalancing.

The KNN classifier was implemented from scratch using a combined Euclidean–Hamming distance function.
Model performance was evaluated for different values of $k$, and the best results were achieved at
$k = 3$, where the classifier obtained the highest accuracy, precision, recall, and F1-score.
This choice of neighborhood size aligns with the patterns observed in the EDA: because the data exhibit
local structure driven by a few strongly discriminative features, a smaller value of $k$ was best able to
capture these local patterns without oversmoothing. The confusion matrix further showed that the classifier
performed well in distinguishing between the two classes, correctly identifying most positive cases while
maintaining a relatively low number of false alarms.

Overall, the study demonstrates that the KNN algorithm—despite its simplicity—can achieve strong predictive
performance on this medical dataset when appropriate preprocessing and distance modeling are applied. The
results highlight the importance of exercise-related features (\texttt{thalach} and \texttt{oldpeak}) as
key indicators of heart disease, while also illustrating the limitations of features such as \texttt{chol}
and \texttt{trestbps} when considered in isolation.

Future work could explore several extensions, including applying cross-validation for more robust estimation
of model performance, optimizing or learning the numeric-to-categorical distance weighting, and comparing
KNN with more advanced models such as logistic regression, decision trees, SVMs, or ensemble methods.
Such approaches may further improve classification accuracy and yield additional insights into the
relationships among features in cardiac risk prediction.